% Physical unit 027; accepted Session 03 after fresh 400-dpi authority replay.
\GordanUnitStart{027}{printed page 22}
\GordanTopFolio{22}
\begingroup\small
punctum $C$. Cujus utriusque lineae extensio intra $A$ et $C$ cum eadem sit; ab $A$ ad $B$ existunt duo aequa itinera, quorum alterum ipsa est linea geodetica I, alterum duabus partibus conficitur, scilicet linea geodetica II intra $A$ et $C$ et linea geodetica I intra $C$ et $B$.

Quodsi superius iter intra $A$ et $B$ linea omnium brevissima esset, nimirum etiam posterius linea esset omnium brevissima, quod fieri nequit, quia posterius iter in puncto $C$ non est linea geodetica.

II. \gq{Ea linea geodetica $l_0$ primi aut tertii generis, quae ducta ab $A$ ad $B$ circulum $-l_0$ bis secat (altero quidem loco in $C$), non est linea omnium brevissima intra $A$ et $B$.}

Cujus rei demonstrandae ratio priori similis est.
Namque puncto $\Pi$ ab $A$ in altera linea geodetica, quae tangit circulum $l_0$, promota ad $C$, tum in priori a $C$ ad $B$; hoc totum iter extensioni $AB$ lineae geodeticae I aequum est.

Est igitur, quod posuimus, superiori argumentandi ratione demonstratum.

\begin{center}\S\ 6.\end{center}

Linea, tangens circulum $l_1$ in puncto $A$, dividit zonam, circulis $l_1$ et $-l_1$ inclusam, in duas partes, quarum alteram nominare placet \gq{$J$}, intra lineam geodeticam, alteram \gq{$E$}, extra lineam geodeticam sitam.

Quare, vi, quae in theorematis modo propositis inest, percognita, perspicuum est; lineam omnium brevissimam intra $A$ et $B(l\psi)$, $B$ sito in parte \gq{$J$}, primi aut tertii generis esse neque habere intra $A$ et $B$ punctum conversionis; $B$ sito in parte \gq{$E$}, secundi aut quarti generis esse et habere intra $A$ et $B$ punctum conversionis; si
\[
 \psi>\omega
\]
primi aut secundi generis esse, si
\[
 \psi<\omega
\]
tertii aut quarti generis esse.

Quaecunque igitur zonae pars punctum $B$ continet, quo in genere linearum geodeticarum ponenda sit linea omnium brevissima intra $A$ et $B$, definiri potest; ac quidem ea linea geodetica, quae omnium brevissima est, circulum $-l_0$ non secat.

\begin{center}\S\ 7.\end{center}

Quae eo spectant, ut, si linea geodetica pars est extensionis periodi, ceteris lineis geodeticis intra $A$ et $B$ sit brevior;
vel \gq{Linea geodetica $l_0$ hac conditione omnium intra $A$ et $B$ brevissima est linea, si angulorum $\varphi$, ad haec puncta pertinentium, differentia non excedit ultra $180^\circ$.}
\endgroup
